% !PRINTED 471
\weberpage{471}
If we now set in (13)

\[
\varDelta = (g\,C^3)^{\frac{1}{7}}, \quad f = h\,C^2\,(g\,C^3)^{-\frac{4}{7}},
\]

then, with the use of (14), the expression (13) becomes

\[
g^{b+2c}\,h^a\,C^{\frac{m}{14}}\,g^{-\frac{m}{7}},
\]

and in the numerator and denominator the factor $C^{\frac{m}{14}}
g^{-\frac{m}{7}}$ can be cancelled, so that everything is expressed
rationally in terms of $g$ and $h$.

Just like $g, h$, the function $u$ also depends only on the ratio of
the variables $x_1 : x_2 : x_3$, and it follows easily that every
other function with the same property which, like $u$, admits the
substitutions of the group $G_{24}$, is representable rationally in
terms of $u, g, h$. For such a function can first of all, by the
general theorems of \S.\,43, be represented as a rational function of
$u$ and the invariants, and indeed in only one way as an integral
function of $u$ which does not exceed the $6$th degree (Bd.\,I,
\S.\,142). The coefficients in this representation are invariants,
and since $u$ depends only on the ratios, the coefficients too can
depend only on the ratios, and are therefore representable rationally
in terms of $g, h$.

\begin{center}
\S.\ 126.

\gesperrt{Reduction of the general resolvent of the seventh degree to
the special one.}
\end{center}

In the preceding section we became acquainted with two forms of the
resolvent of degree $7$ of the form problem, of which the one, the
special, holds for the case that $f = 0$, that is, for the case that
the point sought lies on the fundamental curve, and the general for
the case that it has an arbitrary position.

The quantities $u, g, h$, when they refer to the point $(x)$, we
shall denote by

\weq{1}{u_x = \dfrac{\varDelta^2 z}{C}, \quad g_x = \dfrac{\varDelta^7}{C^3}, \quad h_x = \dfrac{f\varDelta^4}{C^2}}

The general resolvent shall then be denoted by

\weq{2}{R(u_x, g_x, h_x) = R_x = 0}

and the special one arises from it when
%CONT>>
